\documentclass[11pt,reqno]{amsart}
\usepackage[margin=1.15in]{geometry}
\usepackage[expansion=false]{microtype}
\usepackage{xcolor}
\usepackage[colorlinks=true,linkcolor=blue!55!black,citecolor=blue!55!black,urlcolor=blue!55!black]{hyperref}
\hypersetup{pdftitle={Additive representations by powerful integers},pdfauthor={Roy van Rijn}}
\newtheorem{theorem}{Theorem}
\newtheorem{lemma}[theorem]{Lemma}
\newtheorem{proposition}[theorem]{Proposition}
\DeclareMathOperator{\rad}{rad}
\newcommand{\eps}{\varepsilon}
\newcommand{\Z}{\mathbb Z}
\newcommand{\Q}{\mathbb Q}
\newcommand{\C}{\mathbb C}
\newcommand{\Rel}{\ll_{\eps}}
\title[Additive representations by powerful integers]{Additive representations by powerful integers}
\author{Roy van Rijn}
\subjclass[2020]{11D45, 11D72, 11N25, 11G50}
\keywords{powerful numbers, rational reconstruction, bipartite graphs, Campana points}
\begin{document}
\begin{abstract}
For fixed \(k\ge2\), let \(r_k(N)\) count ordered representations of \(N\)
as a sum of two positive \(k\)-full integers. We prove
\(r_2(N)\ll_\eps N^{3/10+\eps}\) and
\(r_k(N)\ll_{k,\eps}N^{1/(k+1)+\eps}\) for \(k\ge3\).
The proof bounds the number of common neighbors in bipartite graphs formed
by omitting one coordinate of the canonical factorization.
The graph bounds also apply to coprime differences of size at least
\(X^{2/k}\) between \(k\)-full integers bounded by \(X\), for \(k\ge3\).
An application to primitive \(k\)-full triples \(a+b=c\le X\) gives an
upper bound of order \(X^{2/k-5/(3k^2)+O(k^{-3})+\eps}\).
For \(k=3\) we obtain the explicit exponent \(21/40+\eps\).
\end{abstract}
\maketitle

\section{Introduction}
A positive integer \(m\) is \(k\)-full if \(p^k\mid m\) for every prime
\(p\mid m\). Write \(r_k(N)\) for the number of ordered pairs of positive
\(k\)-full integers with sum \(N\).
Erd\H{o}s asked whether \(r_2(N)\ll_\eps N^\eps\) for every \(\eps>0\).
Blomer and Sch\"obel \cite[Corollary~2]{BS} proved the uniform bound
\(r_k(N)\ll_{k,\eps} N^{2/(2k+1)+\eps}\). For \(k=2\), the exponent \(2/5\)
also appears in Chojecki's treatment of the pointwise problem \cite[Theorem~1.2]{Cho}.

\begin{theorem}\label{thm:binary}
For every fixed \(k\ge2\) and every \(\eps>0\),
\[
r_k(N)\ll_{k,\eps}N^{\tau_k+\eps},
\qquad
\tau_2=\frac{3}{10},\qquad \tau_k=\frac1{k+1}\quad(k\ge3).
\]
\end{theorem}

The canonical factorization of a \(k\)-full integer is
\[
m=\prod_{j=k}^{2k-1}a_j^j,
\]
where \(a_{k+1},\ldots,a_{2k-1}\) are squarefree and pairwise coprime.
For \(v_p(m)=kq+r\), with \(0\le r<k\), the prime \(p\) occurs in
\(a_k\) to exponent \(q\) when \(r=0\). If \(r>0\), it occurs to exponent
one in \(a_{k+r}\) and to exponent \(q-1\) in \(a_k\). In particular,
\(a_k\) may have a prime factor in common with some \(a_j\) for \(j>k\).

For a chosen \(e\), keep all the coordinates except \(a_e\) on each side.
These two tuples are vertices of a bipartite graph, and a representation
gives an edge. For \(k=2\) and \(e=2\), a pair of cube factors \(b_1,b_2\)
sharing a neighbor satisfies
\[
\left(\frac{a_1c_2}{a_2c_1}\right)^2\equiv \frac{b_2^3}{b_1^3}\pmod N
\]
in the case of coprime summands. If the omitted coordinates are short,
each modular root specifies only a few reduced fractions. Factoring
those fractions and using the exact representation equations recovers
the common neighbor. We use this argument with two choices of \(e\) whose
bounds are effective in complementary ranges. The same graphs count
differences if the modulus is large enough relative to the summands.

\begin{theorem}\label{thm:difference}
Fix \(k\ge3\) and \(\eps>0\). For \(X\ge2\) and an integer
\(n\ge X^{2/k}\), the number of ordered coprime positive \(k\)-full pairs
\(u,v\le X\) with \(u-v=n\) is
\[
                         O_{k,\eps}(X^{1/(k+1)+\eps}).
\]
\end{theorem}

Let \(N_k(X)\) be the number of ordered triples of positive \(k\)-full
integers \(a+b=c\le X\) with \(\gcd(a,b,c)=1\). These are the primitive
positive Campana points for the equal-multiplicity orbifold
\[
\left(\mathbb P^1,\ (1-1/k)([0]+[1]+[\infty])\right).
\]
The identification sends \((a,b,c)\) to \((a:c)\), whose height is \(c\).
Browning and Verzobio \cite[Theorem~1.1]{BV} prove a diagonal saving
\(1/k^2+O(k^{-5/2})\) over \(2/k\) for sufficiently large \(k\).
The binary sum and difference bounds let us fix any summand when
all three are large. A uniform projective curve bound handles the
remaining coordinates.

\begin{theorem}\label{thm:campana}
For each fixed \(k\ge4\) and every \(\eps>0\),
\[
 N_k(X)\ll_{k,\eps}
 X^{\theta_k+\eps},
 \qquad \theta_k=\frac{6k^2-k-4}{(k+1)(3k^2-k-2)}
                 =\frac2k-\frac{5}{3k^2}+O(k^{-3}).
\]
For \(k=3\) we have \(N_3(X)\ll_\eps X^{21/40+\eps}\).
\end{theorem}

\section{Recovering common neighbors}
We work with signed canonical products with coprime absolute values,
each at most \(X\), in
\begin{equation}\label{eq:coeff}
 \kappa_1\prod_{j=k}^{2k-1}a_j^j+
 \kappa_2\prod_{j=k}^{2k-1}c_j^j=n>0.
\end{equation}
The nonzero integers \(\kappa_1,\kappa_2\) are fixed, and the estimates
are uniform in them. Every factor in either term is a unit modulo
\(n\), and \((a_i,c_j)=1\) for all \(i,j\). For sums we take \(X=n\).

\begin{lemma}\label{lem:edge}
For fixed \(e\ge2\), nonzero integers \(P,Q\) and a positive integer \(n\),
the equation \(P a^e+Q c^e=n\) has \(O_{e,\eps}(X^\eps)\) positive
solutions with \(|P|a^e,|Q|c^e\le X\), uniformly in \(P,Q,n\).
\end{lemma}
\begin{proof}
Assume \(e=2\) and order the terms so that \(P>0\). If \(-PQ\) is a
square, then the two integer factors \(Pa\pm c\sqrt{-PQ}\) have
product \(Pn\) and determine \(a,c\). Otherwise, \(Pa+c\sqrt{-PQ}\)
is an algebraic integer of norm \(Pn\) in the quadratic field
\(\Q(\sqrt{-PQ})\). Prime-ideal factorization gives at most \(\tau(Pn)\)
ideals of this norm. In an imaginary quadratic field each principal
ideal has at most six generators of the specified norm. In a real
quadratic field all the displayed generators lie between \(1\) and
\(2X\). The ratio of two generators of the same ideal has norm one
and is a unit. Any such ratio \(w>1\) has integer trace \(w+1/w\ge3\),
which gives \(O(\log(2X))\) generators per ideal in this range.
Since \(P\le X\) and \(n\le2X\), these bounds give the claim.

For \(e\ge3\) and an irreducible form, apply Evertse
\cite[Theorem~1]{Ev} with the real place and the primes dividing \(n\).
The number of \(S\)-unit cosets is at most
\((5\cdot10^6e)^{1+\omega(n)}\). At the fixed positive right-hand
side, homogeneity permits at most one positive pair in each coset.

In the reducible case write \(P X^e+QY^e=G(X,Y)H(X,Y)\) with coprime
integer homogeneous factors of degrees \(d,f>0\). Such a factorization
exists by Gauss's lemma and the fact that the linear factors over \(\C\)
are distinct. Fix the nonzero values \(u=G(a,c)\), \(v=H(a,c)\).
There are at most \(2\tau(n)\) choices because \(uv=n\).
The ratio \(t=a/c\) satisfies
\[
                 v^dG(t,1)^f-u^fH(t,1)^d=0.
\]
This polynomial is nonzero and has degree at most \(df\):
\(G(t,1)\) and \(H(t,1)\) are coprime and nonconstant since
\(Y\) does not divide the original form. Each \(t\) then gives
at most one positive \(c\) by homogeneity.
\end{proof}

Let \(A_j,C_j\) be the lower endpoints of positive dyadic boxes on
\(a_j,c_j\), respectively. Define
\[
 M_1=\prod_{j\ne e}A_j,\qquad
 M_2=\prod_{j\ne e}C_j,\qquad H_e=A_e C_e .
\]

\begin{proposition}\label{prop:graph}
In the graph omitting exponent \(e\), the number of representations
in a box is bounded, with a constant depending only on \(k,\eps\), by
\begin{equation}\label{eq:graph}
 X^\eps\left(M_1+M_2+
     (1+H_e/\sqrt n)\min(M_1\sqrt{M_2},M_2\sqrt{M_1})\right).
\end{equation}
\end{proposition}
\begin{proof}
Consider two distinct vertices on the first side sharing one on the
second. Including the coefficients, their fixed products give
\[
       P_i a_i^e+Q c_i^e=n\qquad (i=1,2).
\]
Here \(a_i,c_i\) refer just to the two omitted coordinates in the
\(i\)-th equation. Take \(a_i=g\alpha_i\) and \(c_i=h\gamma_i\)
by removing each common divisor. The fraction
\(u/v=\alpha_1\gamma_2/(\alpha_2\gamma_1)\) is reduced by the two
gcd reductions and \((a_i,c_i)=1\). Eliminating \(Q\) gives
\begin{equation}\label{eq:recover}
 (u/v)^e\equiv P_2/P_1\pmod n,\qquad
 g^e(P_1u^e-P_2v^e)=n(\gamma_2^e-\gamma_1^e).
\end{equation}
There are \(O_k(e^{\omega(n)})\) unit roots of the first congruence
by the prime-power unit groups and the Chinese remainder theorem.
Within one modular class of fractions with \(u,v\le4H_e\), distinct
fractions have cross difference a nonzero multiple of \(n\).
The corresponding vectors in the first quadrant have lengths
\(O(H_e)\) and angles separated by \(\gg n/H_e^2\).
Thus each class contains \(O(1+H_e^2/n)\) fractions.

A fraction \(u/v\) gives at most \(\tau(u)\tau(v)\) choices of the
four reduced factors. The coefficient of \(g^e\) in \eqref{eq:recover}
is nonzero for an actual common neighbor. If it vanished, then
\(\gamma_1=\gamma_2=1\), so the complete first summands would be
equal, and uniqueness of the canonical factorization would give
the same first vertex. We recover \(g\) and then \(Qh^e\) from either
original equation. Dividing \(Qh^e\) by its fixed signed coefficient
gives a canonical product, since \(h\mid c_i\). It recovers the common
vertex and \(h\).

The codegree bound on both sides is therefore
\(K\ll_{k,\eps}X^\eps(1+H_e^2/n)\).
A bipartite graph with side sizes \(m_1,m_2\) has at most
\(m_1+m_2+\sqrt K\min(m_1\sqrt{m_2},m_2\sqrt{m_1})\) edges under
this bound. Indeed the sum of \(\binom{\deg v}{2}\) over the
second side is at most \(K\binom{m_1}{2}\), and Cauchy--Schwarz
gives one estimate. Interchanging the sides gives the other.
Lemma~\ref{lem:edge} bounds the multiplicity of each edge.
Using \(\sqrt{1+H_e^2/n}\le 1+H_e/\sqrt n\) proves
\eqref{eq:graph}. Counting all possible vertex pairs also gives \(O_{k,\eps}(X^\eps M_1M_2)\).
\end{proof}

\section{Optimizing the boxes}
First take \(k\ge3\) and \(n\ge X^{2/k}\) in \eqref{eq:coeff}.
Let \(x,y,z\) denote, respectively, the
logarithmic dyadic sizes in base \(X\) of the coordinates of
exponents \(k,k+1\) and the sum of the remaining logarithmic sizes
on one side. Write \(x',y',z'\) for the other side, and put
\(t=1/(k+1)\). The size constraint gives
\(kx+(k+1)y+(k+2)z\le1\), and likewise on the other side.
Fixing one complete summand determines the other. We may therefore
assume \(x+y+z\ge t\) and \(x'+y'+z'\ge t\), which imply
\(x\ge z\) and \(x'\ge z'\).

Set \(p=y+z\), \(q=x+z\) and define \(p',q'\) symmetrically.
The constraints give
\begin{equation}\label{eq:complement}
                  p,p'\le t,\qquad q\le 2(t-p),\quad q'\le2(t-p').
\end{equation}
For example, \(2t(kx+(k+1)y+(k+2)z)-(q+2p)
=(k-1)t(x-z)\ge0\). Put \(w=\log n/(2\log X)\ge1/k\).
If \(p+p'\le t\), counting all vertex pairs in the \(e=k\) graph
suffices. Otherwise \(k(x+x')+(k+1)(p+p')\le2\) gives
\(x+x'<1/k\), so the \(e=k\) graph has a bounded height factor.
This graph gives the desired bound unless both \(p+p'/2>t\)
and \(p'+p/2>t\). In that case \(p,p'>t/2\) and \eqref{eq:complement}
gives \(q,q'<t\) and \(q+q'/2<t\). Thus the \(e=k+1\) graph
suffices if \(y+y'\le w\).

If \(y+y'>w\), the exponent of the square-root term in the
\(e=k+1\) graph, including its height factor, is
\[
 \min(q+q'/2,q'+q/2)+y+y'-w
   =s+s'-\tfrac12\max(q,q')-w,
\]
where \(s=x+y+z\) and \(s'=x'+y'+z'\). The size constraints give
\(s\le t(1+x-z)\le t(1+q)\) and similarly on the other side.
The last displayed exponent is therefore at most
\(2t+(2t-\tfrac12)\max(q,q')-w\le t\), since \(t\le1/4\)
and \(w\ge1/k\ge t\).
The number of dyadic boxes is \(O_k((\log(2X))^{2k})\).
Absorbing this factor in \(X^\eps\) proves
Theorem~\ref{thm:difference} and the coprime coefficient bound
for the positive sign with \(X=n\).

For \(k=2\), write the summands as
\(\kappa_1 a^2 b^3\) and \(\kappa_2 c^2 d^3\), with scales
\(A,B,C,D\). In addition to \(AC\) and \(BD\) from
Lemma~\ref{lem:edge}, and the two graph bounds, we use the
mixed-coordinate bound
\begin{equation}\label{eq:mixed}
 \ll_\eps n^\eps\min\bigl(AD\max(B,C)^{1/3},\,
                            BC\max(A,D)^{1/3}\bigr).
\end{equation}
For fixed \(a,d\) the equation in \(b,c\) is \(P b^3+Qc^2=n\).
It defines an absolutely irreducible plane cubic, since
\((n-PX^3)/Q\) has simple zeros and is not a square in \(\C(X)\).
Bombieri--Pila \cite[Theorem~5]{BP} in a square of side
\(2\max(B,C)\) gives the first bound. Fixing \(b,c\) gives the
second.

Let \(\alpha,\beta,\gamma,\delta\) be the logarithms to base \(n\)
of \(A,B,C,D\), and set \(S=\alpha+\gamma,\ W=\beta+\delta\).
The size constraints are
\[
 2\alpha+3\beta\le1,\qquad 2\gamma+3\delta\le1.
\]
The fixed-pair estimates suffice unless \(S,W\ge 3/10\).
Then \(S\le11/20\) and \(W\le7/15<1/2\), so the graph on the
square-root coordinates has a bounded height factor.
If a cube-coordinate exponent is at least \(3/10\), take
\(\beta\ge3/10\) by symmetry, so \(\alpha\le1/20\).
When \(\gamma<3/10\), the graph on the square-root coordinates
suffices, since \(\alpha+\gamma/2<1/5\).
When \(\gamma\ge3/10\), use \eqref{eq:mixed}.
If \(\gamma\ge\beta\), its first exponent is at most
\(1/20+(1-\gamma)/3\le17/60\). If \(\beta\ge\gamma\), it is
at most \((1-3\beta)/2+\beta/3+2/15\le17/60\).

It remains to take \(\beta,\delta<3/10\).
If \(S>1/2\), the graph on the cube coordinates suffices, since
the exponent of its square-root term with the height factor is at
most \(3W/4+S-1/2\le S/2\le11/40\).
If \(S\le1/2\) and a square-root exponent is at least \(3/10\),
take \(\alpha\ge3/10\) by symmetry. Then \(\beta\le2/15\)
and the graph on the cube coordinates again suffices, since
\(\beta+\delta/2<17/60\).
In the remaining case every exponent is below \(3/10\) and both
height factors are bounded. For both graph estimates to have
exponent at least \(3/10\), one must have \(S,W\ge2/5\).
The constraint \(2S+3W\le2\) forces \(S=W=2/5\).
Equality in both graph estimates requires
\(\alpha=\beta=\gamma=\delta=1/5\).
This gives the exponent \(3/10\) in every box.

To pass from coprime summands to arbitrary \(m_1+m_2=N\), take
\(h=(m_1,m_2)\mid N\), \(n=N/h\).
A prime occurring with exponent \(1,\ldots,k-1\) in \(m_i/h\)
divides \(h\), with exponent at least \(k\) in \(h\).
Let \(\kappa_i\) be the product of the full powers of all such
primes in \(m_i/h\). Then \(\kappa_i\mid h\) and
\((m_i/h)/\kappa_i\) is \(k\)-full, so its canonical factorization
puts the quotients in \eqref{eq:coeff}. There are at most
\(\tau(N)^3\) choices of \(h,\kappa_1,\kappa_2\).
The box estimates are uniform in these coefficients. The divisor
function and the \(O_k((\log(2n))^{2k})\) dyadic boxes cost only a
subpower factor. This completes the proof of Theorem~\ref{thm:binary}.

\newpage
\section{Primitive triples and Campana points}
To prove Theorem~\ref{thm:campana}, fix a common canonical dyadic box
for the three summands. Let \(s_i\) be the sum of logarithmic
coordinate sizes to base \(X\) on side \(i\), and \(u_i\) the sum
excluding the coordinate of exponent \(k\). Write
\[
 S=s_1+s_2+s_3,\qquad U=u_1+u_2+u_3,\qquad v_i=s_i-u_i .
\]
The size constraints are \(ks_i+u_i\le1\), which imply \(U\le3-kS\).

Fix the non-\(k\)-power coordinates on all three sides in
\(O_k(X^U)\) ways. The remaining coordinates lie on
\(\lambda x^k+\mu y^k=\nu z^k\) with fixed positive coefficients.
The projective curve is smooth, because the partial derivatives
cannot all vanish at a projective point, and is absolutely
irreducible by B\'ezout. A primitive solution of \(a+b=c\) is
pairwise coprime, so its positive root-coordinate triple is
primitive and injects into the rational points of the curve.
In the rectangular box, Heath-Brown's coefficient-uniform bound
\cite[Theorem~3]{HBcurves} gives the exponent
\[
 U+\frac{v_1+v_2+v_3-\max_i v_i}{k}
    \le U+\frac{2(S-U)}{3k}.
\]
Here the largest monomial size is the \(k\)-th power of the largest
coordinate bound, because the form is diagonal. The same theorem
in an enclosing cube gives the exponent \(U+2/k^2\). When \(k=3\),
Heath-Brown's uniform diagonal-cubic theorem
\cite[Theorem~1.3]{HBcubic} gives the exponent \(U\).

Let \(\Delta\) be the minimum of the available counting exponents
in a box, so its contribution is \(O_{k,\eps}(X^{\Delta+\eps})\).
First count triples all of whose entries are at least \(X^{2/k}\).
Fixing \(c\) and using Theorem~\ref{thm:binary}, or fixing \(a\) or
\(b\) and using Theorem~\ref{thm:difference}, gives
\(\Delta\le s_i+t\) for all \(i\), where \(t=1/(k+1)\).
Thus \(S\ge3(\Delta-t)\). For \(k\ge4\), set
\(L=k-2/3-2/(3k)>0\). The rectangular curve bound gives
\[
 \Delta\le U+\frac{2(S-U)}{3k}
     \le 3-\frac2k-L S
     \le 3-\frac2k-3L(\Delta-t).
\]
Solving yields \(\Delta\le\theta_k\). For \(k=3\) we use
\(\Delta\le U\le3-3S\le3-9(\Delta-1/4)\), so
\(\Delta\le21/40\).

For the remaining triples, by symmetry suppose \(a<X^{2/k}\).
Its size constraint strengthens to \(ks_1+u_1\le2/k\).
Any two summands determine the third, so both
\(\Delta\le s_1+s_2\) and \(\Delta\le s_1+s_3\) hold. In particular
\(S\ge2\Delta-s_1\). For \(k\ge4\), the enclosing-cube bound gives
\[
 \Delta\le U+\frac2{k^2}
       \le 2+\frac2k-kS+\frac2{k^2}
       \le 2+\frac4k+\frac2{k^2}-2k\Delta .
\]
The resulting exponent \((2+4/k+2/k^2)/(2k+1)\) is smaller than
\(2/(k+1)<\theta_k\) for \(k\ge4\).
For \(k=3\), use \(\Delta\le U\) in the same calculation to give
\(\Delta\le10/21<21/40\). In each range the number of boxes is
absorbed in \(X^\eps\).

\paragraph{AI assistance and review.}
ChatGPT (OpenAI) assisted with mathematical arguments, literature checks,
exposition, proof review and verification code.

\begin{thebibliography}{9}

\bibitem{BS}
V.~Blomer and A.~Sch\"obel,
\emph{Twins of powerful numbers},
Functiones et Approximatio Commentarii Mathematici \textbf{49}
(2013), 349--356.
\href{https://doi.org/10.7169/facm/2013.49.2.12}{doi:10.7169/facm/2013.49.2.12}.

\bibitem{BP}
E.~Bombieri and J.~Pila,
\emph{The number of integral points on arcs and ovals},
Duke Math. J. \textbf{59} (1989), 337--357.

\bibitem{BV}
T.~Browning and M.~Verzobio,
\emph{Sums of three powerful numbers},
2026, \href{https://arxiv.org/abs/2608.24512v1}{arXiv:2608.24512v1}.

\bibitem{Cho}
P.~Chojecki,
\emph{Erd\H{o}s Problem 943 on sums of two squarefull numbers},
18 March 2026, \url{https://www.ulam.ai/research/erdos943.pdf}.

\bibitem{Ev}
J.-H.~Evertse,
\emph{The number of solutions of the Thue--Mahler equation},
J. Reine Angew. Math. \textbf{482} (1997), 121--149.

\bibitem{HBcurves}
D.~R.~Heath-Brown,
\emph{The density of rational points on curves and surfaces},
Ann. of Math. \textbf{155} (2002), 553--598.

\bibitem{HBcubic}
D.~R.~Heath-Brown,
\emph{Counting square-full solutions to \(x+y=z\)},
2026, \href{https://arxiv.org/abs/2601.07817v1}{arXiv:2601.07817v1}.

\end{thebibliography}
\end{document}
